\subsection{LLM-Based Requirements Engineering: From Task Automation to Multi-Stage and End-to-End Workflows}
\label{sec:rw-llm-re}

Initial applications of LLMs to requirements engineering have largely focused on automating individual RE activities. Görer and Aydemir developed a GPT-based system for generating requirements elicitation interview scripts for training purposes. Their approach uses a dedicated requirements-elicitation knowledge base together with prompt chaining to construct complete interview scenarios~\cite{gorer2024gpt}.

Moving from training materials to direct interview automation, LLMREI uses an LLM-based chatbot to autonomously conduct requirements elicitation interviews. It evaluates zero-shot and least-to-most prompting and demonstrates that an LLM interviewer can dynamically adapt its questions to stakeholder responses while eliciting a substantial portion of the target requirements~\cite{korn2025llmrei}.

Other studies focus on transforming or analyzing requirements-related artifacts. Almeida et al.\ compare GPT-4 and DeepSeek-V3 for deriving functional and non-functional requirements from transcribed stakeholder interviews~\cite{almeida2025elicitation}. Binkhonain and Alfayez investigate prompt-based LLMs for software requirements classification using zero-shot, few-shot, persona, and chain-of-thought prompting across multiple classification tasks~\cite{binkhonain2025prompts}. For regulatory requirements, Abualhaija et al.\ propose XTRAREG, which combines LLMs with retrieval-augmented generation to extract privacy requirements from predefined legal sources, using the GDPR as a case study~\cite{abualhaija2025llm}.

The capability of LLMs to produce early software requirements has also been compared with human expertise. Hymel and Johnson evaluate LLM-generated and human-generated requirements documents under controlled conditions, showing that LLM-generated requirements can achieve strong alignment with stakeholder ideas while substantially reducing generation time and cost~\cite{hymel2025analysis}.

A related direction applies requirements-engineering principles to human--LLM interaction rather than directly automating a software RE workflow. ROPE frames prompt construction as a requirements-oriented activity in which users explicitly specify and iteratively refine requirements for an LLM's output. Its results indicate that training users to formulate accurate and complete requirements can improve downstream LLM performance~\cite{ma2025rope}.

Prior work has therefore largely focused on isolated RE tasks, while recent systems increasingly coordinate multiple RE activities within broader workflows. In this review, \emph{multi-stage RE} refers to support for multiple RE activities or stages, whereas \emph{end-to-end RE} refers to broader coverage of the requirements-development process from initial stakeholder or requirements input toward downstream requirements artifacts.

MARE proposes an end-to-end multi-agent RE framework containing Stakeholder, Collector, Modeler, Checker, and Documenter agents. The framework iteratively performs elicitation, requirements modeling, verification, and specification, while a shared workspace allows agents to exchange intermediate requirements artifacts~\cite{mare2024}.

iReDev similarly supports end-to-end requirements development across elicitation, analysis, specification, and validation through six knowledge-driven agents---interviewer, end-user, deployer, analyst, archivist, and reviewer. An artifact pool provides event-driven communication among agents, while a human-in-the-loop mechanism incorporates feedback from real-world roles into requirements development~\cite{jin2025iredev}.

VibeIntent also provides an end-to-end multi-stage workflow, but is specifically oriented toward progressive stakeholder-intent refinement for Agile development. Rather than directly generating a final artifact from an initial project description, it progressively transforms stakeholder intent through product vision, elicitation agenda, stakeholder interview, requirements specification, product backlog, and backlog validation artifacts~\cite{vibeintent2026}.

QUARE represents a different type of multi-stage RE system. Rather than supporting the broader end-to-end requirements-development process, it focuses on requirements quality analysis, cross-quality conflict negotiation, KAOS-based integration, and standards compliance verification~\cite{cheng2026quare}.

Overall, LLM-based RE has progressed from isolated task automation toward structured multi-stage and end-to-end workflows. However, although several frameworks include stakeholder or end-user roles, their primary contributions concern RE-task decomposition, artifact processing, knowledge integration, workflow orchestration, or quality analysis. The systematic construction and evaluation of the stakeholder representations themselves remains a comparatively separate research concern. A detailed comparison of these studies is given in Table~\ref{tab:apx-re-studies} of Appendix~\ref{apx:re-studies}.
